% Klausur 1 (29.07.2026): nur Aufgabenstellungen.
% Gemeinsame Präambel für alle LaTeX-Dateien.
\documentclass[a4paper,11pt]{article}
\usepackage[utf8]{inputenc}
\usepackage[T1]{fontenc}
\usepackage[ngerman]{babel}
\usepackage{amsmath,amssymb,amsthm}
\usepackage{enumitem}
\usepackage{booktabs}
\setcounter{MaxMatrixCols}{30}

\theoremstyle{definition}
\newtheorem{definition}{Definition}[section]
\newtheorem{satz}[definition]{Satz}
\newtheorem{lemma}[definition]{Lemma}
\newtheorem{korollar}[definition]{Korollar}
\newtheorem{beispiel}[definition]{Beispiel}
\newtheorem{bemerkung}[definition]{Bemerkung}
\newtheorem{algorithmus}[definition]{Algorithmus}
\newtheorem{aufgabe}{Aufgabe}[section]
\newenvironment{loesung}{\par\noindent\textbf{Lösung.}\ }{\par}

% Kurzbefehle
\newcommand{\N}{\mathbb{N}}
\newcommand{\Z}{\mathbb{Z}}
\newcommand{\Q}{\mathbb{Q}}
\newcommand{\R}{\mathbb{R}}
\newcommand{\Zm}[1]{\mathbb{Z}_{#1}}
\newcommand{\ggT}{\operatorname{ggT}}
\newcommand{\kgV}{\operatorname{kgV}}

\usepackage[a4paper,margin=2.2cm]{geometry}
\usepackage{fancyhdr}
\pagestyle{fancy}\fancyhf{}
\lhead{Diskrete Mathematik -- Klausur 29.07.2026}\rhead{Seite \thepage}
\setlength{\parindent}{0pt}\setlength{\parskip}{4pt}
\renewcommand{\theaufgabe}{\arabic{aufgabe}}

\begin{document}

\begin{center}
{\LARGE\bfseries Klausur Diskrete Mathematik}\\[4pt]
{\large 29.07.2026}
\end{center}

\vspace{6pt}
\begin{tabular}{@{}p{0.5\textwidth}p{0.45\textwidth}@{}}
Name: \hrulefill & Matrikelnr.: \hrulefill
\end{tabular}

\vspace{10pt}

\begin{aufgabe}
Alice will die Nachricht $w = 1511$ (GO) verschlüsseln und als verschlüsselte Nachricht $c$ (Codewort)
an Bob senden. Sie verabreden als Modul die Zahl $m = 43 \cdot 57 = 2451$ zu wählen und den
öffentlichen Schlüssel $e = 221$ zu verwenden.
\begin{enumerate}[label=\alph*)]
  \item Schreiben Sie $e$ im Binärsystem.
  \item Mit welcher Formel verschlüsselt Alice das Wort $w$? Berechnen Sie das Codewort $c$.
  \item Bob besitzt den privaten Schlüssel $d$. Wie (Formel) entschlüsselt Bob das Codewort $c$?
  \item Berechnen Sie $\varphi(m)$.
  \item Oscar erhält Kenntnis von $m$ und $e$. Welche Formeln muss er nutzen, um den privaten
        Schlüssel $d$ zu berechnen? Berechnen Sie $d$.
\end{enumerate}
\end{aufgabe}

\begin{aufgabe}
Wie viele positive ganze Zahlen $\leq 720$ sind weder durch $3$, $4$ oder $5$ teilbar?
Verwenden Sie die Symbole $A_i$, $\alpha_i$ aus der Vorlesung. (Inklusion -- Exklusion)
\end{aufgabe}

\begin{aufgabe}
Bestimmen Sie mit Hilfe des Euklidischen Algorithmus
\begin{enumerate}[label=\alph*)]
  \item den größten gemeinsamen Teiler von $m = 17$ und $n = 15$,
  \item eine ganzzahlige Lösung der Gleichung $17x - 15y = 1$.
\end{enumerate}
\end{aufgabe}

\begin{aufgabe}
Gegeben ist das System simultaner Kongruenzen (Chinesischer Restsatz)
\begin{align*}
  z &\equiv 8 \pmod{17}\\
  z &\equiv 5 \pmod{15}
\end{align*}
\begin{enumerate}[label=\alph*)]
  \item Bestimmen Sie eine ganze Zahl $z$ mit Hilfe des Ergebnisses aus Aufgabe 3.
  \item Bestimmen Sie sämtliche Lösungen dieses Systems, insbesondere die kleinste positive Lösung.
\end{enumerate}
\end{aufgabe}

\begin{aufgabe}
Die Parameter eines Modulo-11-Codes sind $[10, 8, 3]_{11}$.
Vorgelegt ist das Klartextwort $a = 1\,1\,5\,0\,0\,0\,0\,5$.
\begin{enumerate}[label=\alph*)]
  \item Wie lautet die Kontrollmatrix $H$?
  \item Berechnen Sie die Kontrollziffern $c_9$, $c_{10}$ und damit das zugehörige Codewort $c$.
\end{enumerate}
Empfangen wurden die Wörter
\begin{align*}
  r_1 &= 1\,1\,5\,0\,0\,0\,0\,5\,1\,1\\
  r_2 &= 1\,1\,5\,0\,0\,0\,0\,7\,3\,3
\end{align*}
\begin{enumerate}[label=\alph*), start=3]
  \item Berechnen Sie die Syndrome $S(r_1)$, $S(r_2)$.
  \item Welches der beiden Wörter $r_1$, $r_2$ ist nicht dekodierbar und warum nicht?
  \item Dekodieren Sie das andere Wort und geben Sie das Codewort $c$ an.
  \item Zeigen Sie, dass $c$ tatsächlich ein Codewort ist.
\end{enumerate}
\end{aufgabe}

\vfill
\textbf{Punktzahl:} \rule{3cm}{0.4pt} \qquad \textbf{Note:} \rule{3cm}{0.4pt}

\medskip
Bewertungsschema (Punkte/Note):\\[2pt]
{\small
\begin{tabular}{*{10}{c}}
\toprule
29/1.0 & 28/1.3 & 26/1.7 & 25/2.0 & 24/2.3 & 22/2.7 & 21/3.0 & 19/3.3 & 17/3.7 & 15/4.0\\
\bottomrule
\end{tabular}}

\end{document}
